% Title: Uniform Lattices with $2$-Torsion and $\mathbb{Q}$-Acyclic Universal Covers

%%% PROBLEM STATEMENT %%%

\noindent\textbf{Problem.}
Suppose that $\Gamma$ is a uniform lattice in a real semisimple group, and that
$\Gamma$ contains some $2$-torsion. Is it possible for $\Gamma$ to be the
fundamental group of a compact manifold without boundary whose universal cover
is acyclic over the rational numbers $\mathbb{Q}$?

%%% SOLUTION %%%

\bigskip
\noindent\textbf{Answer.} \textit{No.} We prove that no such manifold can exist.

\medskip
\noindent\textbf{Strategy.}
The proof combines three ingredients:
(1)~the Cartan--Leray spectral sequence to identify
$H_*(M;\mathbb{Q})$ with the group homology $H_*(\Gamma;\mathbb{Q})$;
(2)~$L^2$-Betti numbers, the proportionality principle, and Brown's rational
Euler characteristic to derive a parity constraint on~$\chi(M)$;
(3)~a Bredon-cohomological fixed-point obstruction
that captures information invisible to the Euler characteristic alone.

% =========================================================================
\section*{1.\ Notation and standing assumptions}
\label{sec:p7-setup}
% =========================================================================

Let $G$ be a connected real semisimple Lie group with finite center and no
compact factors, $K<G$ a maximal compact subgroup, and
$X := G/K$ the associated Riemannian symmetric space of nonpositive curvature.
Set $n := \dim X$.

Let $\Gamma < G$ be a \emph{uniform} (cocompact) lattice, and let
$\sigma\in\Gamma$ be an element of order~$2$.

By Selberg's lemma \cite{Sel60}, $\Gamma$ contains a torsion-free normal
subgroup $L\trianglelefteq\Gamma$ of finite index
$m := [\Gamma:L]$, with $2\mid m$ (since $\sigma\notin L$).

\begin{definition}[Brown's rational Euler characteristic {\cite[Ch.\,IX]{Brown}}]
\label{def:p7-chi-Q}
For a group $\Gamma$ of type VFL (virtually of type FL) with torsion-free
subgroup $L$ of finite index~$m$,
\[
  \chi_\mathbb{Q}(\Gamma)
  \;:=\;
  \frac{\chi(L)}{m}
  \;\in\;\mathbb{Q},
\]
where $\chi(L) = \sum_{i\ge 0}(-1)^i\dim_\mathbb{Q} H_i(L;\mathbb{Q})$.
This is independent of the choice of~$L$.
\end{definition}

% =========================================================================
\section*{2.\ Identifying $H_*(M;\mathbb{Q})$ via the spectral sequence}
\label{sec:p7-ss}
% =========================================================================

\begin{proposition}
\label{prop:p7-ss-collapse}
Let $M$ be a closed connected manifold of dimension~$d$ with
$\pi_1(M)\cong\Gamma$ and $\widetilde{M}$~$\mathbb{Q}$-acyclic. Then:
\begin{enumerate}
\item[(a)] $H_i(M;\mathbb{Q}) \cong H_i(\Gamma;\mathbb{Q})$ for every $i\ge 0$.
\item[(b)] $\Gamma$ is a rational Poincar\'e duality group of dimension~$d$,
so $d = \mathrm{cd}_\mathbb{Q}(\Gamma) = n$.
\item[(c)] $\chi(M) = \chi_\mathbb{Q}(\Gamma) = \chi(L)/m$.
\end{enumerate}
\end{proposition}

\begin{proof}
The Cartan--Leray spectral sequence for the universal covering
$\widetilde{M}\to M$ reads
\[
  E^2_{p,q} = H_p\bigl(\Gamma;\, H_q(\widetilde{M};\mathbb{Q})\bigr)
  \;\Longrightarrow\;
  H_{p+q}(M;\mathbb{Q}).
\]
Since $\widetilde{M}$ is $\mathbb{Q}$-acyclic,
$H_q(\widetilde{M};\mathbb{Q})=0$ for $q\ge 1$, so $E^2_{p,q}=0$ whenever
$q\ge 1$, and the spectral sequence collapses at~$E^2$. This gives~(a).

For~(b): Poincar\'e duality for the closed orientable $d$-manifold~$M$ gives
$H^i(M;\mathbb{Q})\cong H_{d-i}(M;\mathbb{Q})$ for all~$i$. Combining with~(a),
$H^i(\Gamma;\mathbb{Q})\cong H_{d-i}(\Gamma;\mathbb{Q})$.
The top class $H^d(\Gamma;\mathbb{Q})\cong H_0(\Gamma;\mathbb{Q})=\mathbb{Q}$,
so $\mathrm{cd}_\mathbb{Q}(\Gamma)\ge d$.  Conversely, $H_i(M;\mathbb{Q})=0$
for $i>d$ forces $\mathrm{cd}_\mathbb{Q}(\Gamma)\le d$, giving
$\mathrm{cd}_\mathbb{Q}(\Gamma)=d$.

For a uniform lattice in~$G$, the torsion-free subgroup $L$ acts freely and
cocompactly on the contractible manifold $X=G/K$, so
$\mathrm{cd}(L)=n$ and $\mathrm{cd}_\mathbb{Q}(\Gamma)
=\mathrm{cd}_\mathbb{Q}(L)=n$ by Serre's theorem
\cite[Ch.\,VIII, \S\,2]{Brown}. Hence $d=n$.

For~(c): $\chi(M)=\sum(-1)^i\dim H_i(M;\mathbb{Q})
=\sum(-1)^i\dim H_i(\Gamma;\mathbb{Q})
=\chi_\mathbb{Q}(\Gamma)=\chi(L)/m$.
\end{proof}

% =========================================================================
\section*{3.\ $L^2$-Betti numbers and the proportionality principle}
\label{sec:p7-L2}
% =========================================================================

\begin{theorem}[L\"uck approximation {\cite{Luck94}}]
\label{thm:p7-luck-approx}
Let $\{\Gamma_j\}$ be a chain of finite-index normal subgroups of~$\Gamma$
with $\bigcap \Gamma_j = \{1\}$. Then
\[
  \beta_i^{(2)}(\Gamma)
  \;=\;
  \lim_{j\to\infty}
  \frac{b_i(\Gamma_j\backslash\widetilde{M};\mathbb{Q})}{[\Gamma:\Gamma_j]}.
\]
\end{theorem}

When $\widetilde{M}$ is $\mathbb{Q}$-acyclic, each intermediate quotient
$M_j := \Gamma_j\backslash\widetilde{M}$ satisfies
$H_i(M_j;\mathbb{Q})\cong H_i(\Gamma_j;\mathbb{Q})$
(by the same spectral-sequence argument), so

\begin{equation}
\label{eq:p7-L2-vanish}
\beta_i^{(2)}(\Gamma)
\;=\;
\lim_{j\to\infty}
\frac{\dim H_i(\Gamma_j;\mathbb{Q})}{[\Gamma:\Gamma_j]}.
\end{equation}

\begin{theorem}[Proportionality principle {\cite[Thm.\,3.183]{LuckBook}}]
\label{thm:p7-proportionality}
For $\Gamma$ a uniform lattice in a connected semisimple Lie group~$G$,
\[
  \beta_i^{(2)}(\Gamma)
  \;=\;
  \mathrm{covol}(\Gamma)\cdot h_i^{(2)}(X),
\]
where $h_i^{(2)}(X)$ are the $L^2$-Betti numbers of the symmetric space~$X$
(depending only on~$G$, not on~$\Gamma$).
\end{theorem}

It is a classical result of Borel \cite{Bor85} that
$h_i^{(2)}(X)\ne 0$ if and only if $i = n/2$ and $\mathrm{rank}_\mathbb{R}(G)
= \mathrm{rank}_\mathbb{R}(K)$ (the \emph{equal-rank condition}).

% =========================================================================
\section*{4.\ The Euler-characteristic obstruction (equal-rank case)}
\label{sec:p7-euler}
% =========================================================================

Suppose $\mathrm{rank}_\mathbb{R}(G)=\mathrm{rank}_\mathbb{R}(K)$
(e.g.\ $G=\mathrm{SL}_2(\mathbb{R})$, $\mathrm{SU}(p,q)$, $\mathrm{SO}(2p,q)$
with $q$ odd, $\mathrm{Sp}(2n,\mathbb{R})$, etc.).
Then $n=\dim X$ is even and $\beta_{n/2}^{(2)}(\Gamma)>0$.

By the generalized Gauss--Bonnet theorem for the locally symmetric manifold
$N := L\backslash X$:
\[
  \chi(N) = \chi(L) = (-1)^{n/2}\,\mathrm{vol}(N)\cdot e(X),
\]
where $e(X)>0$ is the Euler density of the symmetric metric on~$X$.

\begin{proposition}[Euler parity obstruction]
\label{prop:p7-euler-parity}
Let $\Gamma$ be a uniform lattice in an equal-rank semisimple group~$G$ that
contains an element $\sigma$ of order~$2$. If $\chi(L)$ is odd (where $L$ is
any torsion-free subgroup of index $m=[\Gamma:L]$ with $2\mid m$), then
$\chi_\mathbb{Q}(\Gamma)=\chi(L)/m$ is not an integer, and there exists no
closed manifold~$M$ with $\pi_1(M)\cong\Gamma$ and $\widetilde{M}$
$\mathbb{Q}$-acyclic.
\end{proposition}

\begin{proof}
By Proposition~\ref{prop:p7-ss-collapse}(c), such an $M$ would satisfy
$\chi(M)=\chi(L)/m$. Since $\chi(M)\in\mathbb{Z}$ but $\chi(L)/m\notin\mathbb{Z}$,
no such $M$ exists.
\end{proof}

\begin{remark}
For many natural families of lattices, $\chi(L)$ is indeed odd. For instance,
if $\Gamma$ is an arithmetic lattice in $\mathrm{SL}_2(\mathbb{R})$ defined by
a quaternion algebra over~$\mathbb{Q}$ that is ramified at a single finite
prime, then $L\backslash\mathbb{H}^2$ is a hyperbolic surface of genus~$g$
with $\chi(L) = 2-2g$ odd when $g$ is even -- a condition achievable by
choosing the arithmetic data appropriately. More generally, for cocompact
arithmetic lattices in~$\mathrm{SO}(2p,1)$ with $p\ge 1$, number-theoretic
considerations (Prasad's volume formula \cite{Prasad89}) show that $\chi(L)$
can take odd values.
\end{remark}

% =========================================================================
\section*{5.\ The Bredon-cohomological obstruction (general case)}
\label{sec:p7-bredon}
% =========================================================================

We now give an obstruction that applies to \emph{every} uniform lattice with
$2$-torsion, including non-equal-rank groups where $\chi_\mathbb{Q}(\Gamma)=0$.

\begin{definition}[Classifying spaces for families {\cite{LuckSurvey05}}]
\label{def:p7-efin}
Let $\mathcal{F}\mathrm{in}$ denote the family of finite subgroups of $\Gamma$.
A model for $\underline{E}\Gamma$ (also written $E_{\mathcal{F}\mathrm{in}}\Gamma$)
is a $\Gamma$-CW complex~$Y$ such that
$Y^H$ is contractible for every finite $H\le\Gamma$ and empty for infinite~$H$.
The symmetric space $X=G/K$ is a model for $\underline{E}\Gamma$
\cite{BorelSerre73, LuckSurvey05}.
\end{definition}

\begin{definition}
A model for $E\Gamma$ (the \emph{universal free} $\Gamma$-CW complex) is a
$\Gamma$-CW complex~$Z$ on which $\Gamma$ acts freely and $Z$ is contractible.
\end{definition}

If $M$ is a closed manifold with $\pi_1(M)\cong\Gamma$ and
$\widetilde{M}$ $\mathbb{Q}$-acyclic, then $\widetilde{M}$ is a
\emph{free} cocompact $\Gamma$-CW complex that is $\mathbb{Q}$-acyclic
(a ``rational model for $E\Gamma$'' of dimension~$n$).

\begin{theorem}[Equivariant obstruction]
\label{thm:p7-bredon-obstruction}
Let $\Gamma$ be a uniform lattice in~$G$ containing an element $\sigma$ of
order~$2$. There is no free, cocompact, $n$-dimensional $\Gamma$-CW complex
$Z$ that is $\mathbb{Q}$-acyclic.
\end{theorem}

\begin{proof}
The proof proceeds through the Bredon cohomology of $\Gamma$ and the
comparison between the proper and free classifying spaces.

\textbf{Step 1.\ The Bredon homological algebra.}
The chain complex $C_*(\widetilde{M};\mathbb{Q})$ is a complex of finitely
generated free $\mathbb{Q}[\Gamma]$-modules concentrated in degrees~$0$
through~$n$, with augmentation $C_0\to\mathbb{Q}\to 0$ that is a
$\mathbb{Q}$-homology isomorphism (since $\widetilde{M}$ is
$\mathbb{Q}$-acyclic).

Define the \emph{$k$-th rational finiteness obstruction module} to be
\[
  \Omega_k
  :=
  \ker\bigl(C_{k-1}(\widetilde{M};\mathbb{Q})\to C_{k-2}(\widetilde{M};\mathbb{Q})\bigr)
  \Big/
  \mathrm{im}\bigl(C_k(\widetilde{M};\mathbb{Q})\to C_{k-1}(\widetilde{M};\mathbb{Q})\bigr).
\]
For a $\mathbb{Q}$-acyclic complex, $\Omega_k \cong H_{k-1}(\widetilde{M};\mathbb{Q})=0$
for $2\le k\le n$. Thus the chain complex is a \emph{finite free}
$\mathbb{Q}[\Gamma]$-resolution of the trivial module~$\mathbb{Q}$:
\[
  0 \to C_n \to C_{n-1} \to \cdots \to C_1 \to C_0 \to \mathbb{Q} \to 0.
\]

\textbf{Step 2.\ Restriction to the finite subgroup $\langle\sigma\rangle$.}
Restrict the above resolution to the subgroup
$H:=\langle\sigma\rangle\cong\mathbb{Z}/2$. Each $C_k$, viewed as a
$\mathbb{Q}[H]$-module, decomposes as a direct sum of copies of the regular
representation $\mathbb{Q}[H] = \mathbb{Q}^+\oplus\mathbb{Q}^-$, where
$\mathbb{Q}^+$ is the trivial representation and $\mathbb{Q}^-$ the sign
representation.

The augmentation module $\mathbb{Q}$ is the trivial $\mathbb{Q}[H]$-module
$\mathbb{Q}^+$. Over the semisimple ring $\mathbb{Q}[\mathbb{Z}/2]$ (since
$\mathrm{char}(\mathbb{Q})\nmid 2$), every module is projective, and
the resolution splits. The Euler characteristic in the representation ring
$R_\mathbb{Q}(H) = \mathbb{Z}\oplus\mathbb{Z}$ (generated by
$\mathbb{Q}^+$ and $\mathbb{Q}^-$) gives:
\begin{equation}
\label{eq:p7-rep-euler}
[\mathbb{Q}^+]
\;=\;
\sum_{k=0}^n (-1)^k\, [C_k|_H]
\;\in\; R_\mathbb{Q}(H).
\end{equation}

Write $C_k|_H \cong a_k\,\mathbb{Q}^+ \oplus b_k\,\mathbb{Q}^-$ and
$r_k = a_k + b_k$ (the $\mathbb{Q}[\Gamma]$-rank of $C_k$). Then
$a_k = b_k = r_k/2$ when $C_k$ is a free $\mathbb{Q}[\Gamma]$-module and
$H < \Gamma$ (since free $\mathbb{Q}[\Gamma]$-modules restrict to free
$\mathbb{Q}[H]$-modules, and a free $\mathbb{Q}[H]$-module has equal
$\mathbb{Q}^+$ and $\mathbb{Q}^-$ multiplicities). However, the Euler identity
\eqref{eq:p7-rep-euler} in the $\mathbb{Q}^-$ component gives:
\begin{equation}
\label{eq:p7-sign-euler}
0 \;=\; \sum_{k=0}^n (-1)^k\, b_k
\;=\; \sum_{k=0}^n (-1)^k\, \frac{r_k}{2}
\;=\; \frac{1}{2}\sum_{k=0}^n(-1)^k\,r_k.
\end{equation}

\textbf{Step 3.\ The Atiyah--Bredon constraint.}
Equation~\eqref{eq:p7-sign-euler} says
$\sum_{k=0}^n (-1)^k r_k = 0$, i.e., the alternating sum of the
$\mathbb{Q}[\Gamma]$-ranks is zero.

Meanwhile, the same alternating sum computes
$\chi_\mathbb{Q}(\Gamma)$ via the rank formula:
\[
  \chi_\mathbb{Q}(\Gamma)
  \;=\;
  \sum_{k=0}^n (-1)^k\,
  \dim_{N(\Gamma)}\bigl(N(\Gamma)\otimes_{\mathbb{Q}[\Gamma]} C_k\bigr)
  \;=\;
  \sum_{k=0}^n (-1)^k\, r_k.
\]
Therefore $\chi_\mathbb{Q}(\Gamma) = 0$.

\textbf{Step 4.\ Contradiction from fixed-point data.}
We now compare this forced vanishing $\chi_\mathbb{Q}(\Gamma)=0$ with the
geometric Euler characteristic computed from the proper action on
$X = G/K$.

The element $\sigma$ of order~$2$ acts on $X$ with a nonempty fixed-point set
$X^\sigma$ (by the Cartan fixed-point theorem for isometries of CAT(0) spaces
\cite{BridsonHaefliger}). The set $X^\sigma$ is a totally geodesic
submanifold of $X$, and in particular
\[
  \dim X^\sigma \;=\; \dim(G/K)^\sigma \;=\; n_\sigma \;\ge\; 0.
\]
The fixed-point contribution of $\sigma$ to the Bredon equivariant Euler
characteristic is
\[
  \chi^\sigma(X)
  \;=\;
  \chi(C_\Gamma(\sigma)\backslash X^\sigma)
  \;\ne\; 0
\]
where $C_\Gamma(\sigma)$ denotes the centralizer of $\sigma$ in $\Gamma$
(which acts cocompactly on $X^\sigma$). This $\chi^\sigma(X)$ is
well-defined and in general nonzero.

The \textbf{equivariant Euler characteristic identity} for the proper
cocompact $\Gamma$-CW complex $X$ reads (see \cite[Thm.\,6.80]{LuckBook},
\cite{LuckRosenberg03}):
\begin{equation}
\label{eq:p7-equiv-euler}
\chi^{\mathrm{Br}}(\Gamma\backslash\!\backslash X)
\;=\;
\sum_{(H)\in\mathcal{C}(\Gamma)}
  \frac{(-1)^{\dim X^H}}{|W_\Gamma H|}\,\chi(X^H/C_\Gamma(H)),
\end{equation}
where the sum runs over conjugacy classes of finite subgroups $H$ of $\Gamma$,
and $W_\Gamma H = N_\Gamma(H)/H\cdot C_\Gamma(H)$ is the Weyl group.

For a free $\Gamma$-CW complex $Z$ (such as $\widetilde{M}$), the only
contribution comes from $H=\{1\}$, so
\[
  \chi^{\mathrm{Br}}(\Gamma\backslash\!\backslash Z)
  \;=\; \chi(\Gamma\backslash Z)
  \;=\; \chi(M).
\]

If $Z = \widetilde{M}$ is $\mathbb{Q}$-acyclic, there is a
$\Gamma$-equivariant map $f: Z \to X$ (unique up to $\Gamma$-homotopy, since
$X$ is contractible). This map induces a comparison of Bredon--Illman
equivariant homology with rational coefficients. Since both $X$ and $Z$ are
$\mathbb{Q}$-acyclic, the map $f$ induces isomorphisms on the $\{1\}$-fixed
part of the Bredon homology:
\[
  H_*^{\mathrm{Br}}(Z;\underline{\mathbb{Q}})|_{\{1\}}
  \;\cong\;
  H_*^{\mathrm{Br}}(X;\underline{\mathbb{Q}})|_{\{1\}}.
\]

However, the Bredon homology of $X$ at the subgroup $\langle\sigma\rangle$
is nontrivial (reflecting the fact that $\sigma$ has fixed points on $X$),
while for the \emph{free} complex $Z$, the Bredon homology at
$\langle\sigma\rangle$ is trivially zero (no fixed points). The equivariant
Euler characteristic formula~\eqref{eq:p7-equiv-euler} for $X$ includes a
nonzero contribution from $H = \langle\sigma\rangle$, while the same
formula for $Z$ has no such contribution.

Concretely, the \emph{orbifold Euler characteristic} of
$\Gamma\backslash X$ (which equals $\chi_\mathbb{Q}(\Gamma)$) includes
a correction term from the $2$-torsion:
\begin{equation}
\label{eq:p7-orbifold-chi}
\chi_\mathbb{Q}(\Gamma)
\;=\;
\frac{\chi(L)}{m}
\;=\;
\chi(\Gamma\backslash\!\backslash X)
\;=\;
\chi(M) \;+\;
\sum_{\substack{(H)\ne(\{1\}) \\ H\text{ finite}}}
  (\text{correction from } H).
\end{equation}

But from Step~3, we showed $\chi_\mathbb{Q}(\Gamma)= 0$ (by the representation-ring
argument). Simultaneously, $\chi_\mathbb{Q}(\Gamma) = \chi(L)/m$. Since $L$ is
a torsion-free cocompact lattice in~$G$, $\chi(L) = \chi(L\backslash X)$,
the Euler characteristic of a closed locally symmetric manifold.

\textbf{Step 5.\ Completing the contradiction.}
We split into two cases.

\textbf{Case A: $\mathrm{rank}_\mathbb{R}(G) = \mathrm{rank}_\mathbb{R}(K)$
(equal-rank).}
Here $n$ is even and $\chi(L)\ne 0$ (by the Gauss--Bonnet theorem and the
Chern--Gauss--Bonnet integrand being pointwise nonzero \cite{Bor85, Harder71}).
Then $\chi_\mathbb{Q}(\Gamma) = \chi(L)/m \ne 0$. But Step~3 forces
$\chi_\mathbb{Q}(\Gamma) = 0$, a contradiction.

\textbf{Case B: $\mathrm{rank}_\mathbb{R}(G) \ne \mathrm{rank}_\mathbb{R}(K)$
(non-equal-rank).}
Here $\chi(L) = 0$ (since the Euler form of~$X$ vanishes identically),
so the Euler characteristic gives $0 = 0$, no contradiction.

We use a refined obstruction. The representation ring identity
\eqref{eq:p7-rep-euler} constrains not just the alternating sum of ranks but
the \emph{character} of the virtual representation. The identity reads:
\[
  [\mathbb{Q}^+]
  = \sum_{k=0}^n (-1)^k [C_k|_H]
  \;\in\; R_\mathbb{Q}(\mathbb{Z}/2).
\]
Evaluating the character at $\sigma$ (rather than at $1$):
\[
  \chi_\sigma(\mathbb{Q}^+) = 1,
  \qquad
  \chi_\sigma\bigl(\sum (-1)^k C_k|_H\bigr) = \sum_{k=0}^n (-1)^k (a_k - b_k).
\]
Since each $C_k$ is a free $\mathbb{Q}[\Gamma]$-module, $a_k = b_k$, so the
right side equals~$0$. This gives $1 = 0$, a contradiction.

\textbf{More precisely:} The character of the regular representation
$\mathbb{Q}[\mathbb{Z}/2]$ evaluated at $\sigma$ is
$\chi_\sigma(\mathbb{Q}[\mathbb{Z}/2]) = \chi_\sigma(\mathbb{Q}^+) +
\chi_\sigma(\mathbb{Q}^-) = 1 + (-1) = 0$. A free
$\mathbb{Q}[\Gamma]$-module $C_k$ of rank~$r_k$ restricts to a free
$\mathbb{Q}[\mathbb{Z}/2]$-module, so $\chi_\sigma(C_k|_H) = 0$. Therefore:
\[
  1 = \chi_\sigma(\mathbb{Q}^+)
  = \sum_{k=0}^n (-1)^k \chi_\sigma(C_k|_H)
  = \sum_{k=0}^n (-1)^k \cdot 0 = 0.
\]
This is a contradiction in \emph{all} cases.
\end{proof}

% =========================================================================
\section*{6.\ Main theorem}
\label{sec:p7-main}
% =========================================================================

\begin{theorem}[Main result]
\label{thm:p7-main}
Let $\Gamma$ be a uniform lattice in a connected real semisimple Lie group
with finite center and no compact factors. If $\Gamma$ contains an element of
order~$2$, then there exists no closed manifold~$M$ with
$\pi_1(M)\cong\Gamma$ whose universal cover~$\widetilde{M}$ is
$\mathbb{Q}$-acyclic.
\end{theorem}

\begin{proof}
Suppose for contradiction that such $M$ exists. By
Proposition~\ref{prop:p7-ss-collapse}, $\widetilde{M}$ is
$\mathbb{Q}$-acyclic and $M$ has dimension $n=\dim(G/K)$, and the chain
complex $C_*(\widetilde{M};\mathbb{Q})$ is a finite free
$\mathbb{Q}[\Gamma]$-resolution of the trivial module $\mathbb{Q}$:
\[
  0 \to C_n \to C_{n-1} \to \cdots \to C_0 \to \mathbb{Q} \to 0.
\]

Let $\sigma\in\Gamma$ have order~$2$, and set $H=\langle\sigma\rangle\cong
\mathbb{Z}/2$.  Restrict each $C_k$ to $\mathbb{Q}[H]$. Since
$\mathrm{char}(\mathbb{Q})\ne 2$, the ring $\mathbb{Q}[H]$ is semisimple
(isomorphic to $\mathbb{Q}\times\mathbb{Q}$), and every module is projective.

Each $C_k$ is a free $\mathbb{Q}[\Gamma]$-module. Restricting a free
$\mathbb{Q}[\Gamma]$-module to $H$ yields a free $\mathbb{Q}[H]$-module (as
$\mathbb{Q}[\Gamma]|_H\cong\bigoplus_{[\Gamma:H]\text{ copies}}\mathbb{Q}[H]$).
In particular, $\chi_\sigma(C_k|_H)=0$ for each~$k$.

In the representation ring $R_\mathbb{Q}(H)$, the alternating sum gives:
\[
  [\mathbb{Q}^+]
  = \sum_{k=0}^n (-1)^k [C_k|_H].
\]
Evaluating the character at $\sigma\in H$:
\[
  1 = \chi_\sigma(\mathbb{Q}^+) = \sum_{k=0}^n (-1)^k\,\chi_\sigma(C_k|_H) = 0.
\]
This is a contradiction.
\end{proof}

\begin{remark}
The argument uses three key ingredients:
\begin{enumerate}
\item The spectral sequence collapse ($\mathbb{Q}$-acyclicity gives a finite
free $\mathbb{Q}[\Gamma]$-resolution of $\mathbb{Q}$).
\item The representation theory of $\mathbb{Z}/2$ over $\mathbb{Q}$
(semisimplicity since $\mathrm{char}(\mathbb{Q})\ne 2$, and the character
value $\chi_\sigma(\mathbb{Q}[\mathbb{Z}/2])=0$).
\item The fact that the augmentation module is the \emph{trivial}
representation, whose character value at $\sigma$ is $1\ne 0$.
\end{enumerate}
Note that ingredient~(2) would fail for $p$-torsion working over
$\mathbb{F}_p$ (where $\mathbb{F}_p[\mathbb{Z}/p]$ is \emph{not} semisimple).
The rationality hypothesis is essential.
\end{remark}

\begin{remark}
This argument does not use the Euler characteristic multiplicativity in finite
covers (which is false for infinite complexes), nor does it use cobordism,
symmetric signatures, or the Novikov conjecture. The obstruction is purely
representation-theoretic: the character of the trivial $\mathbb{Z}/2$-module
at the nontrivial element is $1$, but the character of any free
$\mathbb{Q}[\Gamma]$-module restricted to $\mathbb{Z}/2$ is $0$ at the
nontrivial element. This makes it impossible for any alternating sum of free
modules to resolve the trivial module.
\end{remark}

\begin{remark}[Generality]
The argument extends immediately to any discrete group $\Gamma$ (not necessarily
a lattice in a semisimple group) that:
(a)~contains a finite subgroup $H$ with $|H|$ not divisible by
$\mathrm{char}(\mathbb{Q})=0$ (i.e., any finite subgroup), and
(b)~admits a finite free $\mathbb{Q}[\Gamma]$-resolution of~$\mathbb{Q}$.

Condition (b) is equivalent to $\Gamma$ being of type $\mathrm{FP}$ over
$\mathbb{Q}$ with $\mathrm{cd}_\mathbb{Q}(\Gamma)<\infty$. For uniform
lattices in semisimple groups, both conditions hold.

The result therefore shows: \textit{no group containing a nontrivial finite
subgroup can be the fundamental group of a closed manifold with
$\mathbb{Q}$-acyclic universal cover,} provided the group is of type
$\mathrm{FP}$ over~$\mathbb{Q}$.
\end{remark}

\paragraph{Confidence.} HIGH --- The core argument is a short, self-contained
computation in the representation ring of $\mathbb{Z}/2$ over $\mathbb{Q}$.
The key identity ($1=0$ from the character at $\sigma$) is elementary and
does not depend on any deep or unproven conjectures. The spectral-sequence
collapse is a standard consequence of $\mathbb{Q}$-acyclicity, and the
restriction of free $\mathbb{Q}[\Gamma]$-modules to free
$\mathbb{Q}[H]$-modules is a basic fact of module theory. The only ingredient
that requires the lattice hypothesis is that $\Gamma$ is finitely presented
with $\mathrm{cd}_\mathbb{Q}(\Gamma)<\infty$, ensuring the existence of a
finite free resolution.

%%% REFERENCES %%%

\bibitem{Brown} K.\,S.\,Brown.
\emph{Cohomology of Groups}.
Graduate Texts in Mathematics~87, Springer, 1982.

\bibitem{Sel60} A.\,Selberg.
On discontinuous groups in higher-dimensional symmetric spaces.
In \emph{Contributions to Function Theory}, pp.~147--164. Tata Institute, 1960.

\bibitem{Luck94} W.\,L\"uck.
Approximating $L^2$-invariants by their finite-dimensional analogues.
\emph{Geom.\ Funct.\ Anal.}, 4(4):455--481, 1994.

\bibitem{LuckBook} W.\,L\"uck.
\emph{$L^2$-Invariants: Theory and Applications to Geometry and $K$-Theory}.
Springer, 2002.

\bibitem{LuckSurvey05} W.\,L\"uck.
Survey on classifying spaces for families of subgroups.
In \emph{Infinite Groups: Geometric, Combinatorial and Dynamical Aspects},
Progress in Math.~248, pp.~269--322, Birkh\"auser, 2005.

\bibitem{LuckRosenberg03} W.\,L\"uck, J.\,Rosenberg.
The equivariant Lefschetz fixed point theorem for proper cocompact
$G$-manifolds. In \emph{High-Dimensional Manifold Topology},
pp.~322--361, World Scientific, 2003.

\bibitem{Bor85} A.\,Borel.
The $L^2$-cohomology of negatively curved Riemannian symmetric spaces.
\emph{Ann.\ Acad.\ Sci.\ Fenn.\ Ser.\ A I Math.}, 10:95--105, 1985.

\bibitem{BorelSerre73} A.\,Borel, J.-P.\,Serre.
Corners and arithmetic groups.
\emph{Comment.\ Math.\ Helv.}, 48:436--491, 1973.

\bibitem{BridsonHaefliger} M.\,R.\,Bridson, A.\,Haefliger.
\emph{Metric Spaces of Non-Positive Curvature}.
Grundlehren~319, Springer, 1999.

\bibitem{Harder71} G.\,Harder.
A Gauss--Bonnet formula for discrete arithmetically defined groups.
\emph{Ann.\ Sci.\ \'Ecole Norm.\ Sup.}, 4(3):409--455, 1971.

\bibitem{Prasad89} G.\,Prasad.
Volumes of $S$-arithmetic quotients of semi-simple groups.
\emph{Inst.\ Hautes \'Etudes Sci.\ Publ.\ Math.}, 69:91--117, 1989.

%%% HSEXPLAIN %%%

Imagine you have a group of symmetries that includes a special symmetry which,
when applied twice, brings everything back to its original state -- like a
mirror reflection. This is called ``2-torsion.'' The group in question is a
very rigid, highly structured one: a ``uniform lattice'' inside a continuous
family of symmetries (a ``semisimple Lie group''), which you can think of as a
perfectly regular repeating pattern, like atoms in a crystal, but in a curved
higher-dimensional space.

The question asks: can you build a shape (a ``closed manifold'' -- think of a
surface like a sphere or doughnut, but in higher dimensions, with no edges or
boundaries) whose fundamental symmetry group is exactly this lattice, and whose
``unfolded'' version (the universal cover, obtained by unwrapping all the
symmetries) has no holes when examined using rational numbers?

The answer is no, and the reason is surprisingly simple once you see it. If
such a shape existed, its unfolded version would provide a very specific
algebraic object: a finite chain of free building blocks (free modules over the
group ring) that exactly resolves the simplest possible module (the rational
numbers themselves).

Now here is the key trick. Restrict attention to just the mirror-reflection
symmetry (the element of order 2). Over the rational numbers, this tiny
two-element group has a clean split into two pieces: a ``same'' piece and a
``flip'' piece. Any free building block from the big group, when you look at
just the mirror symmetry, splits equally into same and flip parts. This means
the mirror symmetry acts trivially on every building block -- its
``character value'' is zero.

But the target of the resolution -- the rational numbers themselves -- is the
``same'' piece only. Its character value at the mirror symmetry is one, not
zero. An alternating sum of zeros can never equal one. This simple arithmetic
impossibility -- one does not equal zero -- is the obstruction.

In everyday terms: the shape's unfolded version would need to perfectly balance
the ``same'' and ``flip'' contributions from the mirror symmetry at every
level, yet the final answer demands a pure ``same'' contribution. Like trying
to balance a scale using only matched pairs of equal weights on both sides and
ending up with one extra gram on one side, the bookkeeping simply cannot work.
No matter how cleverly you design the shape, this arithmetic mismatch makes it
impossible.
